\section{Behavioural Observation and Detection Engineering}
\label{sec:behavioural-detection}

Chapter~5 established the controlled reproduction result and separated the
decisive privilege transition from the accompanying filesystem observations.
The same separation is necessary for detection engineering. A detector may
observe a prerequisite, an action in the exploitation chain, a change to the
target, or the final privileged outcome; these signals occur at different
times and support different conclusions.

The most repeatable host-side effect in the laboratory was the changed content
returned when the selected SUID executable was read through the normal
mounted/VFS path. Its ownership, size, and permission bits remained unchanged,
but its SHA-256 digest differed from the value recorded before exploitation.
This made content integrity a useful detection point even though it did not
identify the proof-of-concept file or explain the internal kernel state that
caused the changed view.

To test that opportunity, a small integrity monitor named MORI was deployed on
VM~271 while it still ran the reconstructed vulnerable kernel. MORI compared a
protected known-good baseline with repeated reads of selected privileged
executables. During a further controlled Copy Fail execution, it reported a
digest mismatch for \code{/usr/bin/su}. The accompanying capture also showed
that the proof of concept obtained a UID~0 shell. The experiment therefore
demonstrated detection of the security-relevant effect, but not prevention of
the exploit.

A separate exact-sample YARA comparison was used to illustrate the difference
between identifying one known proof-of-concept artifact and observing a
host-side consequence. The comparison is deliberately narrow: it evaluates one
rule whose sole condition was the known sample's SHA-256 digest. It is not an
evaluation of YARA as a whole.

This chapter documents the detection model, candidate observables, protected
target baseline, monitor implementation, controlled result, exact-sample
comparison, and operational limitations. It then answers RQ2 and RQ3 within the
evidence actually preserved by the experiment.

\subsection{Detection Model and Evidential Boundary}

The detection question was framed around the privilege boundary rather than
around one exploit filename. The experimentally demonstrated target was
\code{/usr/bin/su}, but the asset being protected was the wider ability to
execute trusted code with elevated credentials. The observed path can be
summarised as:

\begin{lstlisting}
unprivileged local process
    -> vulnerable AF_ALG AEAD path
    -> changed mounted view of a privileged executable
    -> privileged executable launched
    -> UID 0 process
\end{lstlisting}

This representation combines evidence from the vulnerability analysis,
controlled reproduction, and detection phases. It does not mean that every
transition was captured by one telemetry source. In particular, the preserved
detection dataset does not contain a syscall trace, eBPF event stream, or
kernel trace proving the \code{AF\_ALG}, pipe, \code{splice}, and AEAD
operations during the monitored run. Those operations describe the analysed
Copy Fail mechanism and provide candidates for later correlation; the
integrity mismatch and UID~0 outcome are the events directly preserved for
this detection experiment.

The evidence package contains the following principal records:

\begin{table}[H]
    \centering
    \caption{Evidence used in the behavioural detection evaluation.}
    \label{tab:detection-evidence}
    \small
    \begin{tabularx}{\textwidth}{@{}>{\bfseries}p{0.08\textwidth}>{\raggedright\arraybackslash}p{0.31\textwidth}Y@{}}
        \toprule
        ID & Preserved form & Evidential role \\
        \midrule
        00--01 & Text artifacts and screenshots &
        Record the root-owned SUID inventory and which targets were readable
        and executable by \code{labuser}. \\

        02--03 & Baseline and deployment records &
        Preserve the expected target digests and the root-owned operational
        baseline used by the monitor. \\

        04 & Journal excerpt &
        Records MORI's alert for \code{/usr/bin/su}, including the expected
        and observed digests. \\

        07 & Service-status artifact &
        Records the running service, monitored-file count, configured interval,
        and retained alert in the journal. \\

        11--12 & Screenshots &
        Show the alert alongside proof-of-concept execution and, in the latter
        capture, the resulting UID~0 identity confirmation. \\

        YARA & Rule, scan output, and screenshots &
        Compare an exact-hash match for the original sample with a non-match
        after a one-byte, non-executed variation. \\
        \bottomrule
    \end{tabularx}
\end{table}

Screenshot number~08 is intentionally absent from the sequence because the
planned controlled-alert capture was not retained. The numbering is preserved
rather than revised retrospectively. Screenshots~11 and~12 provide useful
concurrent visual context, but the proof-of-concept launch shown in them has no
embedded timestamp. Consequently, the alert timestamp cannot be subtracted
from a verified execution timestamp to produce a measured detection latency.

The public SHA-256 manifests were also rechecked during report preparation.
All listed detector source, service, machine-readable artifacts, rules, and
screenshots used here matched their recorded digests. The two explanatory
\code{README.md} files---the main detection README and the YARA-comparison
README---did not match the digests listed for them. They are therefore treated
as supporting documentation, not as integrity-verified evidence. The
experimental claims below rely on the separately verified artifacts, source,
rule, and captures.

The term \emph{behavioural detection} is used in a host-centric sense in this
chapter: a decision is based on a security-relevant state transition produced
by the tested technique, rather than on the name or exact bytes of the program
that initiated it. MORI observes an effect of the behaviour, not the complete
causal chain. It can report that a protected executable no longer matches its
known-good mounted-view content, but that alert alone cannot attribute the
change specifically to Copy Fail.

\subsection{Observable Stages and Detection Value}

The reproduction exposed several possible observation points. Their defensive
value depends on whether they are direct evidence of exploitation, merely a
precondition, or an outcome that occurs after the security boundary has
already been crossed. \Cref{tab:detection-opportunities} distinguishes these
roles.

\begin{table}[H]
    \centering
    \caption{Candidate observables and their value in the tested attack path.}
    \label{tab:detection-opportunities}
    \small
    \begin{tabularx}{\textwidth}{@{}>{\bfseries\raggedright\arraybackslash}p{0.22\textwidth}>{\raggedright\arraybackslash}p{0.31\textwidth}Y@{}}
        \toprule
        Observable & Experimental status & Detection assessment \\
        \midrule
        Vulnerable kernel and loaded \code{algif\_aead} &
        Both states were recorded before controlled reproduction. &
        Useful exposure context, but neither state proves malicious activity
        and legitimate workloads may require the interface. \\

        AF\_ALG, pipe, \code{splice}, and AEAD sequence &
        Established by the analysed technique, but not directly traced during
        the integrity-monitor run. &
        Potentially useful for selective correlation; not validated as a
        runtime detector in this phase. \\

        Changed privileged-file content through the mounted view &
        Repeatedly observed for \code{/usr/bin/su} and captured by MORI. &
        Strong, sample-independent integrity signal for covered targets, but
        cause-agnostic and detected only after the changed view exists. \\

        Unchanged size, owner, and mode &
        Recorded alongside the changed digest in Chapter~5. &
        Demonstrates that metadata-only monitoring would miss the observed
        content change. \\

        Execution producing UID~0 &
        Directly shown by \code{id} and \code{whoami}. &
        High-confidence confirmation of impact, but too late to be a preventive
        signal in the demonstrated sequence. \\

        Mounted/backing-view divergence and reboot recovery &
        Preserved in the reproduction evidence. &
        Valuable diagnostic context for the transient effect; unsuitable as a
        lightweight continuous detector in the form tested. \\
        \bottomrule
    \end{tabularx}
\end{table}

The module and kernel observations are useful for exposure management. A host
running an affected kernel with \code{algif\_aead} available presents the
tested preconditions, whereas a corrected kernel or an effective module-loading
restriction changes that exposure. These are configuration findings rather
than exploit alerts: the interface may be loaded for legitimate cryptographic
use, and an affected kernel can remain idle indefinitely.

The syscall-level sequence is closer to attacker activity, but individual
operations are not necessarily unique to Copy Fail. A high-fidelity detector
would need to correlate operation type, ordering, process or thread context,
target relationship, and timing while retaining negative controls for benign
AF\_ALG and pipe use. The dataset in this phase does not perform that
correlation. Treating the analysed sequence as already validated telemetry
would therefore overstate the evidence.

The mounted-view digest deviation was the strongest directly tested detection
opportunity. It occurred on the protected side of the interaction, after the
untrusted process had influenced a privileged target but without requiring
knowledge of how the initiating script was named, encoded, or launched. The
same observation appeared in the preliminary investigation, the preserved
reproduction, and the monitored detection test. This supports repeatability on
the same reconstructed guest; it is not a statistical success-rate measurement
across hosts or workloads.

\subsection{Privileged Target Surface and Known-Good Baseline}

The initial target inventory identified 13 root-owned SUID files on the tested
filesystem. From the UID~1001 \code{labuser} account, 12 were recorded as both
readable and executable. The remaining file,
\code{/usr/lib/dbus-1.0/dbus-daemon-launch-helper}, was readable but not
executable in that account context. The experimentally demonstrated target was
\code{/usr/bin/su}; the other reachable SUID files were treated as candidate
privileged targets, not as individually proven Copy Fail targets.

Executables carrying Linux file capabilities were inventoried separately. They
were not added to this first MORI baseline. The distinction matters because the
inventory described a broader privileged-execution surface than the monitor
actually covered. A clean result from MORI therefore means only that the 12
listed SUID paths matched their expected content at the time of a scan.

The baseline stored one SHA-256 digest and one absolute path per monitored
file. For \code{/usr/bin/su}, the expected value was:

\begin{lstlisting}
c74311fe5636b7d7f9a56239fa8adeeab12ba86fe7d41b91afa85bf9bbdae78b
\end{lstlisting}

The complete baseline file had the following digest:

\begin{lstlisting}
c5ebe8afe16326aea860317b153981a43e57f3d462c3e1d03d3d1bbe5d951e74
\end{lstlisting}

Its operational copy was placed at
\code{/etc/copyfail-detector/suid-known-good.tsv}. The preserved deployment
record shows a root-owned directory and a baseline file with mode
\code{-rw-r--r--}, owner \code{root}, group \code{root}, and size
1,016~bytes. An accompanying screenshot records a \code{labuser} modification
attempt being denied. The machine-readable artifact establishes the ownership
and mode; the denied write is screenshot evidence rather than a separately
preserved command transcript.

Root ownership prevents the ordinary test account from silently replacing the
expected values under the recorded permissions, but it is not a complete trust
solution. The baseline must already be correct when generated, legitimate
package updates require controlled regeneration, and an attacker who has
obtained unrestricted root access may be able to alter the baseline, detector,
service definition, or retained logs. MORI's baseline is consequently a
laboratory known-good reference, not an externally attested measurement.

Content hashing was selected because Chapter~5 showed that the target's size,
ownership, and SUID mode remained at their original values while its
mounted-view digest changed. A metadata-only watch would therefore have
reported no relevant difference in that run. Conversely, a digest mismatch
establishes only that the bytes returned by the monitored read differ; it does
not identify the changed offsets, explain their semantics, or prove whether
the persistent backing data was replaced.

\subsection{MORI Integrity Monitor and Service Deployment}

The detector implementation was intentionally small so that its decision could
be traced directly. At start-up, it loads the protected TSV baseline. It then
opens each listed path through the ordinary filesystem interface, reads it in
65,536-byte chunks, calculates SHA-256, and compares the result with the
expected digest. The tested logic can be represented as:

\begin{lstlisting}
baseline = load protected path-to-hash mappings
repeat:
    for each protected path:
        observed = SHA256(read path through normal filesystem view)
        if a new mismatch is observed: emit ALERT
        if a previous mismatch has cleared: emit RECOVERED
    sleep 0.25 seconds
\end{lstlisting}

Reading through the normal filesystem path is central to this experiment. It
places the monitor on the same mounted/VFS view that returned the changed
digest during reproduction. MORI does not inspect the proof-of-concept process
or calculate the proof-of-concept's hash; it repeatedly measures the protected
target content instead.

The set of paths currently in alert state suppresses identical messages on
every scan and permits a later recovery message if a target again matches its
baseline. Missing or unreadable files enter a separate error path. The
implementation does not cryptographically protect its journal output, stop a
changed file from being executed, quarantine a process, or restore target
content. Its decision is a binary equality test against the configured
baseline.

MORI was deployed as a \code{systemd} service with output directed to the
system journal. The service used \code{DynamicUser=yes} and
\code{NoNewPrivileges=yes}, together with a strict read-only system view,
private temporary and device namespaces, kernel and control-group protection,
and \code{RestrictSUIDSGID=yes}. These settings reduce the monitor's own
privilege and writable surface while still allowing it to read the
world-readable protected executables. They do not make the service or journal
tamper-proof against an already privileged attacker.

The status artifact records the service as enabled and active, with the Python
monitor running under PID~3992. Its startup journal entries state:

\begin{lstlisting}
[2026-08-09T15:27:04.563102+00:00] Watching 12 privileged files
[2026-08-09T15:27:04.563133+00:00] Baseline: /etc/copyfail-detector/suid-known-good.tsv
[2026-08-09T15:27:04.563143+00:00] Check interval: 0.25 seconds
\end{lstlisting}

At the later status capture, the service had been active for approximately
1~hour 12~minutes and reported 6.0~MiB of memory, a 6.5~MiB peak, and
52.947~seconds of accumulated CPU time. These values show that the service was
operational for more than an immediate demonstration, but they are a single
status observation rather than a controlled performance benchmark. No disk-I/O
profile, workload comparison, or production-scale target count was measured.

The nominal 0.25-second sleep is likewise not a measured maximum alert delay.
Each scan also requires the time needed to open and hash all listed files, and
the scheduler may add further delay. A sufficiently short-lived mismatch could
occur entirely between observations. Event-driven monitoring or kernel-level
telemetry could reduce that polling limitation, but neither was evaluated by
this detector.

\subsection{Controlled Detection Result}

With the service running on VM~271, the Copy Fail proof of concept was executed
again from the unprivileged \code{labuser} account. Before the monitored
execution, \code{/usr/bin/su} matched the known-good digest. MORI subsequently
emitted the following journal record:

\begin{lstlisting}
[2026-08-09T15:32:31.501808+00:00] ALERT
path=/usr/bin/su
expected=c74311fe5636b7d7f9a56239fa8adeeab12ba86fe7d41b91afa85bf9bbdae78b
observed=44900c631391f0d60eb6d271b8374a08dc1d9be76e403390d27a91ed5f179be9
\end{lstlisting}

The expected value is the original digest recorded before exploitation. The
observed value is the same changed mounted-view digest documented in the
preserved reproduction. This agreement links the detector decision to the
previously characterised host-side effect without requiring the monitor to
recognise the initiating sample.

Screenshots~11 and~12 place the journal alert next to the
\code{labuser} execution terminal. Screenshot~11 shows the proof-of-concept
command returning a \code{\#} prompt; screenshot~12 adds the explicit identity
checks:

\begin{lstlisting}
# id
uid=0(root) gid=1001(labuser) groups=1001(labuser),100(users)
# whoami
root
\end{lstlisting}

As in Chapter~5, the prompt alone is not treated as proof of privilege. The
\code{id} and \code{whoami} output establish the UID~0 outcome. The journal
artifact independently preserves the mismatch event, while the screenshot
provides the concurrent visual relationship between alert and successful
exploitation.

\begin{table}[H]
    \centering
    \caption{Result of the monitored Copy Fail execution.}
    \label{tab:mori-detection-result}
    \small
    \begin{tabularx}{\textwidth}{@{}>{\bfseries\raggedright\arraybackslash}p{0.25\textwidth}>{\raggedright\arraybackslash}p{0.29\textwidth}Y@{}}
        \toprule
        Question & Direct result & Interpretation \\
        \midrule
        Was the protected target changed in the monitored view? &
        MORI reported the documented expected and changed digests for
        \code{/usr/bin/su}. &
        Yes, for the covered target and the tested execution. \\

        Did detection depend on the PoC filename or hash? &
        The detector read only its target baseline and protected paths. &
        No; its decision was independent of the initiating sample's identity. \\

        Was privilege escalation prevented? &
        The adjacent identity checks returned UID~0 and \code{root}. &
        No; the alert was visibility into a successful exploit. \\

        Was detection latency measured? &
        The alert has a journal timestamp, but the visible launch does not. &
        No defensible time-to-alert value can be calculated. \\
        \bottomrule
    \end{tabularx}
\end{table}

No false-positive rate can be derived from this single targeted experiment.
The clean-state captures show that the baseline could be loaded and monitored
without an immediate alert, and the later status record shows no additional
preserved mismatch beyond the \code{/usr/bin/su} event. They do not constitute
a representative benign workload trial. Legitimate package replacement,
administrative maintenance, or any other change to a covered executable would
also produce a digest mismatch until the baseline was reviewed and updated.

\subsection{Exact-Sample YARA Comparison}

The comparison experiment used YARA~4.5.0 on the same
\code{6.8.0-116-generic} environment. Its test rule imported YARA's
\code{hash} module and matched only when the complete scanned file had the
known proof-of-concept SHA-256 digest:

\begin{lstlisting}
d401e7d1c00605749d6c617ace73ab20a762b72e41c2e1590331596e38219a61
\end{lstlisting}

The original 731-byte laboratory sample was scanned without being executed.
At \code{2026-08-09 19:10:18 UTC}, the output recorded:

\begin{lstlisting}
CopyFail_Known_PoC_Exact_Hash /home/labuser/copy_fail_exp.py
yara exit=0
\end{lstlisting}

A private copy was then changed only by appending one trailing newline. Its
size became 732~bytes and its SHA-256 digest became:

\begin{lstlisting}
a567d09b15f6e4440e70c9f2aa8edec8ed59f53301952df05c719aa3911687f9
\end{lstlisting}

The modified file was not executed. At
\code{2026-08-09 19:14:21 UTC}, the exact-hash rule returned no match. The
experiment therefore proves that this particular exact-sample signature
identified the byte-for-byte sample for which it was written and failed after
a trivial byte-level change. It does not prove that the modified file remained
operational, evaded a more general rule, or would evade other static, semantic,
or behavioural detection methods.

This narrow result must not be generalised into a criticism of YARA itself.
YARA rules may use strings, byte patterns, structural properties, modules, and
compound conditions rather than a whole-file digest. The comparison isolates
the limitation of exact-sample identity: it answers whether the scanned bytes
are the one known artifact, not whether the host has experienced the
security-relevant effect associated with Copy Fail.

MORI makes the opposite trade-off. It does not identify the initiating program
at all. A renamed, re-encoded, or otherwise byte-different implementation would
remain outside the detector's concern if it produced no protected-target
change. If it did cause a covered SUID executable to differ from its
known-good mounted-view content for long enough to be scanned, MORI would make
the same comparison regardless of the initiating filename, language, or hash.
That generalisation across initiators is useful, while the loss of causal
specificity remains an operational limitation.

\subsection{Limitations and Answers to RQ2 and RQ3}

The detection result is affirmative but deliberately bounded. The principal
limitations are:

\begin{itemize}
    \item
    only \code{/usr/bin/su} was experimentally demonstrated as a Copy Fail
    target; the other baseline entries are coverage candidates rather than
    individually validated exploit targets;

    \item
    the baseline covered 12 readable and executable root-owned SUID files, not
    file-capability executables, every privileged program, or arbitrary target
    paths;

    \item
    polling can miss an integrity deviation that begins and ends between
    complete scans, and the 0.25-second sleep is not a measured alert latency;

    \item
    a digest mismatch is cause-agnostic and may also result from an authorised
    update, corruption, or another attack technique;

    \item
    the detector alerts after the changed target view exists and did not stop
    the demonstrated UID~0 transition;

    \item
    no syscall or kernel trace directly validated the proposed lower-level
    correlation signals during this monitored run;

    \item
    no representative benign-workload corpus, false-positive rate,
    multi-host trial matrix, or controlled performance benchmark was collected;
    and

    \item
    baseline, detector, and journal trust may be lost after unrestricted root
    compromise unless measurements and alerts are protected or exported
    independently.
\end{itemize}

Within those limits, RQ2 can be answered by separating consistently observed
effects from merely promising telemetry. On the tested VM, the changed
mounted-view content of \code{/usr/bin/su} was the most consistently reproduced
and directly useful system-level detection opportunity. It was visible in the
preliminary investigation, the preserved reproduction, and the monitored
detection run, while size, ownership, and permissions remained unchanged.
The resulting UID~0 identity was a reliable impact confirmation but occurred
too late to support prevention. Vulnerable-kernel and module state supplied
exposure context, whereas the lower-level AF\_ALG and \code{splice}
correlation remained a design candidate rather than a validated detector in
this phase.

RQ3 can also be answered affirmatively for the demonstrated target and
laboratory conditions. MORI detected the exploitation effect without matching
the proof-of-concept filename, programming language, executable hash, or exact
file contents. The one-byte YARA comparison reinforces why that distinction
matters: the exact-hash rule ceased to match after a trivial, non-executed
variation, whereas MORI's decision was anchored to the protected target's
deviation from known-good state.

The answer is not that MORI uniquely identifies Copy Fail. It detects a
security-relevant integrity violation consistent with the reproduced
technique. Attribution would require additional correlation, and timely
prevention requires a control placed earlier in the attack path. The next
chapter therefore moves from observing the changed privileged target to
evaluating a compensating control intended to interrupt the known exploitation
sequence before the privilege boundary is crossed.
